\documentclass{elsart3-1}

% Use the option doublespacing or reviewcopy to obtain double line spacing
% \documentclass[doublespacing]{elsart}

\usepackage{latexsym}
\usepackage{amsmath}
\usepackage{amsfonts}
\usepackage{amssymb}
\usepackage{graphpap}
\usepackage{epsfig}
\usepackage{amscd}
% \usepackage{amsthm}
\usepackage{enumerate}
\usepackage{eucal} % za polepanje ukaza \cal, t.j. kaligrafije

\newcommand{\C}[0]{{\mathbb C}}
\newcommand{\N}[0]{{\mathbb N}}
\newcommand{\R}[0]{{\mathbb R}}
\newcommand{\cK}[0]{{\mathcal K}}
\newcommand{\f}[0]{\varphi}
\newcommand{\pd}[0]{\partial}
\newcommand{\nor}[1]{\| #1 \|}
\newcommand{\mn}[2]{\{ #1\, ;#2 \}}
\newcommand{\sk}[2]{\langle #1 , #2\rangle}
\newcommand{\Sk}[2]{\Big\langle #1 , #2\Big\rangle}

\newcommand{\RR}{{\mathbb{R}}}
\newcommand{\NN}{{\mathbb{N}}}
\newcommand{\ZZ}{{\mathbb{Z}}}
\newcommand{\CC}{{\mathbb{C}}}
\newcommand{\PP}{{\mathbb{P}}}

\newcommand{\eps}{\varepsilon}
\newcommand{\po}{\partial\Omega}



\newcommand{\bp}{\noindent {\it Proof}.\,\,}
\newcommand{\ep}{\hfill$\Box$ \vskip 0.08in}

\newcommand{\meanint}{{\int{\mkern-19mu}-}}
\def\ring{\mathaccent"0017 }

\newcommand{\todo}[1]{\vspace{5 mm}\par \noindent
\marginpar{\textsc{}} \framebox{\begin{minipage}[c]{0.95
\textwidth} \tt #1
\end{minipage}}\vspace{5 mm}\par}


%\usepackage{showkeys}

\usepackage[english,francais]{babel}


%ENVIRONMENTS THEOREMS...
% These are predefined, and follow the numbering system used in the journal!
%English
\newtheorem{theorem}{Theorem}[section]
\newtheorem{lemma}[theorem]{Lemma}
\newtheorem{e-proposition}[theorem]{Proposition}
\newtheorem{corollary}[theorem]{Corollary}
\newtheorem{e-definition}[theorem]{Definition\rm}
\newtheorem{remark}{\it Remark\/}
\newtheorem{example}{\it Example\/}
%French
\newtheorem{theoreme}{Th\'eor\`eme}[section]
\newtheorem{lemme}[theoreme]{Lemme}
\newtheorem{proposition}[theoreme]{Proposition}
\newtheorem{corollaire}[theoreme]{Corollaire}
\newtheorem{definition}[theoreme]{D\'efinition\rm}
\newtheorem{remarque}{\it Remarque}
\newtheorem{exemple}{\it Exemple\/}
\renewcommand{\theequation}{\arabic{equation}}
\setcounter{equation}{0}

%%%%%%%%%%%%%%%%%%%%%%%%%%%%%%%%
%% GUILLEMETS (FRENCH QUOTES) %%
%%%%%%%%%%%%%%%%%%%%%%%%%%%%%%%%
\def\og{\leavevmode\raise.3ex\hbox{$\scriptscriptstyle\langle\!\langle$~}}
\def\fg{\leavevmode\raise.3ex\hbox{~$\!\scriptscriptstyle\,\rangle\!\rangle$}}


\journal{}
\begin{document}
% place in the next line the header (rubrique) chosen for your article,
% if you know it (you can also have 2, format : Header1/Header2
\centerline{}
\begin{frontmatter}

% Title, authors and addresses

% use the thanksref command within \title, \author or \address for footnotes;
% use the ead command for the email address,
% and the form \ead[url] for the home page:
% \title{Title\thanksref{label1}}
% \thanks[label1]{}
% \author{Name\thanksref{label2}}
% \ead{email address}
% \ead[url]{home page}
% \thanks[label2]{}
% \address{Address\thanksref{label3}}
% \thanks[label3]{}
\selectlanguage{english}
\title{Square function and Riesz transform  in non-integer dimensions}


% use optional labels to link authors explicitly to addresses:
% \author[label1,label2]{}
% \address[label1]{}
% \address[label2]{}
% The [label1] can be suppressed if there is only one address for all authors

\selectlanguage{english}
\author[SM]{Svitlana Mayboroda},
\ead{svitlana@math.purdue.edu}
\author[AV]{Alexander Volberg}
\ead{volberg@math.msu.edu, A.Volberg@ed.ac.uk}

\address[SM]{Department of Mathematics, Purdue University, 150 N. University Street, West Lafayette, IN 47907-2067, USA}
\address[AV]{Department of Mathematics, Michigan State University, East Lansing, MI 48824, USA}

% If you know the dates of reception, and acceptation you can put them now;
%  idem the name of the person presenting the Note

\medskip
\begin{center}
%{\small Received *****; accepted after revision +++++\\
%Presented by }
\end{center}



\begin{abstract}
\selectlanguage{english}
Following a recent paper \cite{RdVTolsa} we show that the finiteness of square function associated with the Riesz transforms with respect to Hausdorff measure $H^s$ implies that $s$ is integer. 
% Text of abstract in English
%{\it To cite this article: A. Name1, A. Name2, C. R. Acad. Sci. Paris, Ser. I 340 (2005)}

\vskip 0.5\baselineskip

\selectlanguage{francais}
% Text of abstract in French
%\noindent{\bf R\'esum\'e} \vskip 0.5\baselineskip \noindent
%{\bf La transformation de Riesz et une fonction carr\'ee dans les dimension fractionnelles.}

%On peut modifier l'article recent \cite{RdVTolsa} pour d\'emontrer que la convergence de la fonction carr\'ee associ\'ee aux transformations de Riesz de mesure de Hausdorff $H^s$ implique que $s$ est un nombre entier.
%{\it Pour citer cet article~: A. Name1, A. Name2, C. R. Acad. Sci.
%Paris, Ser. I 340 (2005).
\end{abstract}
\end{frontmatter}

% now the Version franaise abrge, if it exists
% \selectlanguage{francais}
% \section*{Version fran\c{c}aise abr\'eg\'ee}
% Text of your Version franaise abrge here.
% Note you do not need to repeat here equations that you use in the
% main text - for example 'voir (3)' is quite acceptable.





\selectlanguage{english}
% main text
\section{Introduction}
\label{}
%\setcounter{equation}{0}

% The Appendices part is started with the command \appendix;
% appendix sections are then done as normal sections
% \appendix

\font\tretamajka=cmbsy10 at 12pt

For a Borel mesure $\mu$ in $\RR^m$ and $s\in (0,m]$ the $s$-Riesz transform of $\mu$ is defined as
\begin{equation}\label{eq1.1}
R^s \mu(x):=\int\frac{x-y}{|x-y|^{s+1}}\,d\mu(y), \qquad x\not\in {\rm supp}\,\mu,
\end{equation}

\noindent and the truncated Riesz transform is given by
\begin{equation}\label{eq1.2}
R^s_{\eps}\, \mu(x):=\int_{|x-y|>\eps}\frac{x-y}{|x-y|^{s+1}}\,d\mu(y), \quad R^s_{\eps,\eta}\, \mu(x):=R^s_{\eta}\, \mu(x)-R^s_{\eps}\, \mu(x),
\end{equation}

\noindent where $x\in\RR^m,$ $\eta>\eps>0$.

Further, recall that the upper and lower $s$-dimensional densities of $\mu$ at $x$ are given by  
\begin{equation}\label{eq1.2.1}
\theta^{s,\ast}_{\mu}(x):=\limsup_{r\to 0}\frac{\mu(B(x,r))}{r^s} \quad\mbox{and}\quad \theta^{s}_{\mu,\ast}(x):=\liminf_{r\to 0}\frac{\mu(B(x,r))}{r^s},
\end{equation}

\noindent respectively, where $B(x,r)$ is the ball of radius $r>0$ centered at $x\in \RR^m$. 

It has been proved in \cite{Pr1} and \cite{Pr2} that whenever $0\leq s\leq 1$ and $\mu$ is a finite Radon measure with $0<\theta^{s,\ast}_{\mu}(x)<\infty$, for $\mu - {\rm a.e.}\,x\in\RR^m$,
the condition 
\begin{equation}\label{eq1.2.2}
\sup_{\eps>0} |R_{\eps}^s\,\mu(x)|<\infty\qquad \mu - {\rm a.e.}\,x\in\RR^m,
\end{equation}

\noindent implies that $s\in\ZZ$. Moreover, an analogous result has been obtained in \cite{Vi} for all $0\leq s\leq m$ under a stronger assumption that  $0<\theta^{s}_{\mu,\ast}(x)\leq \theta^{s,\ast}_{\mu}(x) <\infty$. However, neither the curvature methods of \cite{Pr1}, \cite{Pr2}, nor the tangent measure techniques in \cite{Vi} could be applied to establish that \eqref{eq1.2.2} implies $s\in\ZZ$ for all $0\leq s\leq m$ assuming only $0<\theta^{s,\ast}_{\mu}(x)<\infty$.

In  \cite{RdVTolsa} the authors proved that the latter statement holds if the condition \eqref{eq1.2.2} is substituted by the existence of  the principal value $\lim_{\eps\to 0} R_{\eps}^s\,\mu(x)$, $\mu - {\rm a.e.}\,x\in\RR^m$. In the present work we refine the techniques of  \cite{RdVTolsa} and establish the following result. 

\begin{theorem}\label{tSqF} Let $\mu$ be a finite Radon measure in $\RR^m$ with 
\begin{equation}\label{eq1.2.3}
0<\theta^{s,\ast}_{\mu}(x)<\infty\quad\mbox{for}\quad \mu - {\rm a.e.}\,x\in\RR^m.
\end{equation}

\noindent Furthermore, assume that  for some $s\in (0,m]$ the square function
\begin{equation}\label{eq1.3}
S^s\mu(x):=\Bigg(\int_0^\infty \left|R^s_{t,2t}\, \mu(x)\right|^2\,\frac{dt}{t}\Bigg)^{1/2}<\infty, \qquad \mu - {\rm a.e.}\,x\in\RR^m.
\end{equation}

\noindent Then $s\in\ZZ$.
\end{theorem}

%Here $R^s_{2^{-k},2^{-k+1}}\, \mu(x)$, $k\in\ZZ$, $\mu\in\RR^m$, is a version of doubly truncated Riesz transform, defined, for example, as follows. Let $\psi^0$ be a non-decreasing radial $C^2$ function with $\chi{(B(0,1/2))}\leq \psi^0\leq \chi(B(0,2))$, and set $\psi_k(z):=\psi^{0}(2^kz)-\psi^0(2^{k+1}z)$, $k\in\ZZ$, $z\in\RR^m$, so that ${\rm supp}\,\psi_k\subset B(0,2^{-j+1})\setminus B(0,2^{-j-2})$ and $\sum_{k\in\ZZ}\psi_k=1$ in $\RR^m\setminus\{0\}$.  Then we set $R^s_{2^{-k},2^{-k+1}}\, \mu(x)$ to be an operator defined analogously to \eqref{eq1.1} with the kernel given by $\psi_k(x-y)\,\frac{x-y}{|x-y|^{s+1}}$. 

In fact, we also prove the following closely related result which is a strengthening of  the main results in \cite{RdVTolsa}.\begin{theorem}\label{tLim} Let $\mu$ be a finite Radon measure in $\RR^m$  satisfying \eqref{eq1.2.3}. Furthermore, assume that  for some $s\in (0,m]$ we have
\begin{equation}\label{eq1.3.1}
\lim_{\eps\to 0}R^s_{\eps,2\eps}\,\mu(x)=0 \qquad \mu - {\rm a.e.}\,x\in\RR^m.
\end{equation}

\noindent Then $s\in\ZZ$.
\end{theorem}


This circle of problems goes back, in particular, to the work of David and Semmes \cite{DS1}, \cite{DS2}, where the authors showed, under certain assumptions on the measure $\mu$, that the $L^2$ boundedness of a large class of  singular integral operators implies that $s$ is an integer and $\mu$ is uniformly rectifiable, that is, the support of $\mu$ contains ``large pieces of Lipschitz graphs" -- see \cite{DS1}, \cite{DS2} for details. The ultimate goal, which seems to be out of reach at the moment, is to prove that a similar conclusion holds purely under the assumption that the Riesz transform is bounded in $L^2$, i.e., that the Riesz transform alone encodes the geometric information about the underlying measure. The achievements in \cite{NTV-Acta}, \cite{VBook}, \cite{NTV-IMRNWeakType} showed that the $L^2$-boundedness of the Riesz transform, suitably interpreted, is almost equivalent to the condition \eqref{eq1.2.2}. However, under the assumption \eqref{eq1.2.2} the problem seems to be just as challenging.
In both cases the question has only been resolved for $s=1$ (\cite{MMV}, \cite{To2}, \cite{To10}), by the methods involving curvature of measures. 

In this vein, we would like to point out that by Khinchin's inequality \eqref{eq1.3} can be viewed {\it almost} as a condition
\begin{equation}\label{eq1.4}
{\mathbb{E}}\,\Bigg|\sum_{k\in\ZZ} \eps_k R^s_{2^{-k},2^{-k+1}}\, \mu(x)\Bigg|<\infty \qquad \mu - {\rm a.e.}\,x\in\RR^m,
\end{equation}

\noindent where $\eps_k$ are independent random variables taking the values $-1$ and $1$ with probability $1/2$ each. Therefore,  in order to guarantee $s\in\ZZ$, it is sufficient to assume only that the singular integrals of the type $\sum_{k=0}^{\infty} \eps_k R^s_{2^{-k},2^{-k+1}}\, \mu(x)$ are uniformly bounded.

Finally, the $s$ dimensional Hausdorff measure $H^s$ of a set $E$ with $0<H^s(E)<\infty$ satisfies the condition \eqref{eq1.2.3}, and hence, the results of Theorems~\ref{tSqF} and \ref{tLim} remain valid in this context.
 






\section{Preliminary estimates}
%\setcounter{equation}{0}

Our proof largely relies on the estimates for a slightly modified version of the Riesz transform that were obtained in \cite{RdVTolsa}. To be precise, let us consider the operator
\begin{equation}\label{eq2.1}
R^{s,\varphi}_{\eps}\, \mu(x):=\int\varphi\left(\frac{|x-y|^2}{\eps^2}\right)\,\frac{x-y}{|x-y|^{s+1}}\,d\mu(y), \qquad \eps>0,
\end{equation}

\noindent where $\varphi=\varphi_\rho$ is a $C^2$ function depending on the parameter $\rho\in(0,1/2)$, to be determined below, with ${\rm supp}\,\varphi\subset [0,1+\rho+2\rho^2]$ and such that
\begin{enumerate}
\item $\varphi(r)=r^{\frac{s+1}{2}}$ for $0\leq r \leq 1$,
\qquad $\varphi(r)=-\frac r\rho+1+\rho+\frac 1\rho$ for $1+\rho^2\leq r\leq 1+\rho^2+\rho$,
\item $|\varphi(r)|\leq C,$ $|\varphi'(r)|\leq 1/\rho$, $|\varphi''(r)|\leq C(\rho)$ for all $r>0$.
\end{enumerate} 

Analogously to \cite{RdVTolsa}, we start with a set 
\begin{eqnarray}\label{eq2.2}
\nonumber F_\delta &:=& \{x\in \RR^m:\,\mu (B(x,r))/r^s\leq 2\theta_\mu^{s,\ast}(x)\,\mbox{ for }\,r\leq r_0,\quad \theta_\mu^{s,\ast}(x)\leq C_0,\\[4pt]
&&\quad \mbox{and } |R^{s,\varphi}_{\eps}\,\mu(x)-R^{s, \varphi}_{2\eps}\,\mu(x)|\leq \delta \,\mbox{ for all }\,0<\eps<\eps_0\},
\end{eqnarray}

\noindent where $0<\delta<1$ and $C_0, r_0,\eps_0$ are some positive constants.  

One can see that both the condition \eqref{eq1.3} and \eqref{eq1.3.1} imply that 
\begin{equation}\label{eq2.4}
\lim_{\eps\to 0}|R^{s,\varphi}_{\eps}\,\mu(x)-R^{s,\varphi}_{2\eps}\,\mu(x)|=0 \qquad \mu - {\rm a.e.}\,x\in\RR^n.
\end{equation}
\noindent Indeed, 
\begin{equation}\label{eq2.3}
R^{s,\varphi}_{\eps}\, \mu(x)=\int\int_{0<t<\frac{|x-y|^2}{\eps^2}}\varphi'(t)\,dt\,\frac{x-y}{|x-y|^{s+1}}\,d\mu(y)=\int_0^{1+\rho+2\rho^2} \varphi'(t)R_{\eps\sqrt t}^{s}\,\mu(x)\,dt, 
\end{equation}

\noindent so that  $\eqref{eq1.3.1}$ directly gives \eqref{eq2.4}. Furthermore,  $\eqref{eq2.3}$ entails that
\begin{equation}\label{eq2.3.1}\nonumber
\left|R^{s,\varphi}_{\eps}\, \mu(x)-R^{s,\varphi}_{2\eps}\, \mu(x)\right|\leq C(\rho) \int_0^{1+\rho+2\rho^2} |R_{\eps\sqrt t,2\eps\sqrt t}^{s}\,\mu(x)|\,dt\leq C(\rho)\left(\int_0^{\eps\sqrt{1+\rho+2\rho^2}}|R_{u,2u}^{s}\,\mu(x)|^2\,\frac{du}{u}\right)^{1/2},
\end{equation}

\noindent where we used the change of variables $u:=\eps\sqrt t$ and H\"older's inequality. Hence, \eqref{eq1.3} also leads to \eqref{eq2.4}.

 Therefore, for sufficiently small $\eps_0$ and $r_0$ and sufficiently large $C_0$ the set $F_\delta$ has $\mu(F_{\delta})>0$. 
Note that $\mu (B(x,r))\leq Mr^s$ for all $x\in F_\delta$, $r>0$ and $M=\max\{2C_0,\mu(\RR^m)/r_0^s\}$. 

Let $\theta^s(x,r)$ denote  the average $s$-dimensional density of the ball $B(x,r)$, $x\in\RR^m$, $r>0$, that is, 
$\theta^s(x,r):=\mu(B(x,r))/r^s$. We start with the following estimates.

\begin{proposition}\cite{RdVTolsa} \label{p2.1}
Assume that for some $C'>0$, $r>0$ and $x_0\in\RR^m$ we have 
$\mu(B(x_0,r))\geq C'r^s$, and denote by $n$ the biggest integer strictly smaller than $s$. Then for a sufficiently small $\rho$ (depending on $s$ only) and any $\tau\in (0,1/20)$ there exists a constant $\omega_0=C(C',M,\rho)\tau^{-s\,\frac {1}{\log_4(1+\rho^2/4)}}$, there exists $\eps\in\left(\frac r\tau, \omega_0\,\frac r\tau\right]$ and a set of points $y_0,...,y_{n+1}\in B(x_0,r)\cap F_{\delta}$ such that 
\begin{equation}\label{eq2.9}
\theta^s(y_0,4\eps)\leq C(\rho)\,\theta^s(y_0,\eps),\qquad \theta^s(y_0,\eps)\geq C'\tau^s/2,
\end{equation}

\noindent and 
\begin{equation}\label{eq2.7}
\sum_{j=1}^{n+1}|R^{s,\varphi}_{\eps}\,\mu(y_j)-R^{s,\varphi}_{\eps}\,\mu(y_0)|+\theta^s(y_0, 3\eps)\frac{r^2}{\eps^2}\geq C(C',M,s)(n+1-s)\,r\,\frac{\theta^s(y_0,\eps)}{\eps}.
\end{equation}
\end{proposition}

\section{The proof of the main result}
%\setcounter{equation}{0}

%\todo{Note on the proof below: we come to \eqref{eq3.8} essentially without any assumptions on $r,\delta$ and $\tau$ other than $r/\delta<\eps_0$. It is at this moment that we choose $\delta$ and $r$ so that the right-hand side of \eqref{eq3.8} is sufficiently small. The powers of $\tau$ are just convenient to "measure" smallness. So it looks as though we actually shrink $r$ and $\delta$. But the point is that $\theta^s(y_0,\eps)$ stays sufficiently big - controlled from below by $\tau^s$, no matter how small is $r$ (and it has nothing to do with $\delta$). And also, of course, it is important that $\eps$ itself stays within desirable limits. I am not sure if we want to try to explain any of this.}
 
We will argue by contradiction. We initially assume that $s\not\in \ZZ$, and then show that the estimate in \eqref{eq2.7} is accompanied by the corresponding bound from above in terms of $\delta$, $r$, $\tau$ and $\eps$. Ultimately, choosing $r$, $\delta$, $\tau$ sufficiently small leads to a contradiction. Observe that in the case $s\in\ZZ$ the lower bound in \eqref{eq2.7} is degenerate, and hence, such an argument could not be constructed.

Set 
\begin{equation}\label{eq3.1}
\delta:=\frac{\tau^{s+2}}{\omega_0}=C(M,\rho)\, \tau^{s+2+s\,\frac {1}{\log_4(1+\rho^2/4)}},
\end{equation}

\noindent where the constant $C(M,\rho)$ is equal to the reciprocal of $C(C',M,\rho)$ from the definition of $\omega_0$ corresponding to $C'=1/2$. Note that $M$ depends on $r_0$ and $C_0$ in the definition of $F_\delta$, however, the choice of $r_0$ and $C_0$ is determined solely by the properties of $\mu$ and can be made independent of $\delta$.

Going further, fix $\eps_0$ and take 
\begin{equation}\label{eq3.2}
r<\eps_0\delta=\eps_0\,\frac{\tau^{s+2}}{\omega_0}=C(M,\rho)\,\eps_0\, \tau^{s+2+s\,\frac {1}{\log_4(1+\rho^2/4)}}
\quad \mbox{such that} \quad \mu(B(x_0,r)\cap F_\delta)\geq r^s/2.\end{equation}

Now that $r$ and $\delta$ are fixed, we invoke Proposition~\ref{p2.1}, find the points $y_0,..., y_{n+1}$  and choose $\eps\in \left(\frac r\tau,\omega_0\,\frac r\tau\right]$ such that \eqref{eq2.9} and \eqref{eq2.7} are satisfied. However,  for every $x,z\in B(x_0, r)\cap F_{\delta}$, $x\in \RR^m$ we have
\begin{equation}\label{eq3.4}
|R^{s,\varphi}_{\eps}\,\mu(x)-R^{s,\varphi}_{\eps}\,\mu(z)|\leq C(\rho)M \delta + C \delta \log \frac {r}{\delta\eps},
\end{equation}

\noindent whenever  $2\eps<\frac r\delta<\eps_0$. Indeed, a direct calculation shows that for $\eta\in\left[\frac{r}{2\delta},\frac{r}{\delta}\right]$
\begin{equation}\label{eq3.6}
|R^{s,\varphi}_{\eta}\,\mu(x)-R^{s,\varphi}_{\eta}\,\mu(z)|\leq C(\rho)\left(r/\delta\right)^{-s-1}\, |z-x| \,\mu(B(x_0, 4r/\delta))\leq C(\rho)M \delta.
\end{equation}

\noindent Then we can choose $\eta\in\left[\frac{r}{2\delta},\frac{r}{\delta}\right]$ such that $\eta=2^k\eps$ for some $k\in\NN$, so that
\begin{eqnarray}\label{eq3.7}
&&|R^{s,\varphi}_{\eps}\,\mu(x)-R^{s,\varphi}_{\eps}\,\mu(z)|\nonumber\\[4pt]\nonumber &&\qquad\leq |R^{s,\varphi}_{\eps}\,\mu(x)-R^{s,\varphi}_{\eta}\,\mu(x)|+
|R^{s,\varphi}_{\eta}\,\mu(x)-R^{s,\varphi}_{\eta}\,\mu(z)|+|R^{s,\varphi}_{\eta}\,\mu(z)-R^{s,\varphi}_{\eps}\,\mu(z)| \\[4pt] &&\qquad
\leq C(\rho)M \delta + C \sup_{x\in F_{\delta}}\sup_{1\leq i\leq k}|R^{s,\varphi}_{2^{i}\eps}\,\mu(x)-R^{s,\varphi}_{2^{i-1}\eps}\,\mu(x)| \log \frac {r}{\delta\eps}\leq C(\rho)M \delta + C \delta \log \frac {r}{\delta\eps}.
\end{eqnarray}

Therefore, \eqref{eq2.7} is complemented by the estimate
\begin{equation}\label{eq3.5}
\sum_{j=1}^{n+1}|R^{s,\varphi}_{\eps}\,\mu(y_j)-R^{s,\varphi}_{\eps}\,\mu(y_0)|+\theta^s(y_0, 3\eps)\frac{r^2}{\eps^2}\leq 
C(M,\rho)\left(\delta+\delta \log \frac{r}{\eps\delta}+ \theta^s(y_0, \eps)\frac{r^2}{\eps^2}\right),\end{equation}

\noindent where we used \eqref{eq3.4} and \eqref{eq2.9}. Now combining \eqref{eq2.7} with \eqref{eq3.5} and dividing both sides by $r/\eps$ we arrive at the estimate
\begin{equation}\label{eq3.8}
\theta^s(y_0,\eps)\leq C(M,s,\rho)\left(\frac{\delta\eps}{r}+\frac{\delta\eps}{r}\log \frac{r}{\eps\delta}+\theta^s(y_0, \eps)\frac{r}{\eps}\right).
\end{equation}

\noindent According to our choice of $\delta$ and $r$, 
\begin{equation}\label{eq3.9}
\frac{\delta\eps}{r}\leq \frac{\tau^{s+2}}{\omega_0}\,\frac{\omega_0}{\tau}=\tau^{s+1}\leq C\tau \,\theta^s(y_0, \eps) \quad \mbox{and}\quad \frac{r}{\eps} \leq \tau.
\end{equation}

 Now \eqref{eq3.8} and \eqref{eq3.9} give the bound
\begin{equation}\label{eq3.10}
\theta^s(y_0,\eps)\leq C(M,s,\rho)\left(\tau +\tau^{1-\alpha}\right)\theta^s(y_0, \eps), \qquad \forall\,\alpha>0,
\end{equation}

\noindent which for $\tau>0$ sufficiently small leads to a contradiction. \ep






\begin{thebibliography}{00}
% please try to use the bibitem system -
% the references should be in alphabetical order of authors' names.
% Articles with a single author first, author will 1 co-author next,
% then author with several co-authors;


% \bibitem{label}
% Text of bibliographic item


\bibitem{DS1} G.\,David, S.\,Semmes, {\it
Singular integrals and rectifiable sets in $R\sp n$: Beyond Lipschitz graphs.}  
Astrisque No. 193 (1991), 152 pp.

\bibitem{DS2} G.\,David, S.\,Semmes, {\it Analysis of and on uniformly rectifiable sets.} Mathematical Surveys and Monographs, 38. American Mathematical Society, Providence, RI, 1993.

\bibitem{MMV} P.\,Mattila, M.\,Melnikov, J.\,Verdera, {\it The Cauchy integral, analytic capacity, and uniform rectifiability.} Ann. of Math. (2) 144 (1996), no. 1, 127--136.

\bibitem{NTV-IMRNWeakType} F.\,Nazarov, S.\,Treil, A.\,Volberg,
{\it Weak type estimates and Cotlar inequalities for Calder\'on-Zygmund operators on nonhomogeneous spaces. }
Internat. Math. Res. Notices 1998, no. 9, 463--487. 

\bibitem{NTV-Acta} F.\,Nazarov, S.\,Treil, A.\,Volberg, {\it The $Tb$-theorem on non-homogeneous spaces.} Acta Math. 190 (2003), no. 2, 151--239. 

\bibitem{Pr1} L.\,Prat, {\it 
Potential theory of signed Riesz kernels: capacity and Hausdorff measure.}
Int. Math. Res. Not. 2004, no. 19, 937--981. 

\bibitem{Pr2} L.\,Prat, {\it Principal values for the signed Riesz kernels of non-integer dimensions.} Preprint, 2006, as cited in \cite{RdVTolsa}.

\bibitem{To2} X.\,Tolsa, {\it Principal values for the Cauchy integral and rectifiability.} Proc. Amer. Math. Soc. 128 (2000), no. 7, 2111--2119. 

\bibitem{To10} X.\,Tolsa, Xavier {\it
Growth estimates for Cauchy integrals of measures and rectifiability.} 
Geom. Funct. Anal. 17 (2007), no. 2, 605--643. 

\bibitem{RdVTolsa} X.\,Tolsa, A.\,Ruiz de Villa, {\it Non existence of principal values of signed Riesz transforms of non integer dimension.} Preprint, 2008. 

\bibitem{Vi} M.\,Vihtil\"a, {\it
The boundedness of Riesz $s$-transforms of measures in $\mathbf R\sp n$.}  
Proc. Amer. Math. Soc. 124 (1996), no. 12, 3797--3804. 

\bibitem{VBook} A.\, Volberg, {\it
Calder\'on-Zygmund capacities and operators on nonhomogeneous spaces. }
CBMS Regional Conference Series in Mathematics, 100. Published for the Conference Board of the Mathematical Sciences, Washington, DC; by the American Mathematical Society, Providence, RI, 2003.

\end{thebibliography}

\bigskip


\end{document}